\chapter{More on local structure}

This chapter serves for the second half of the commutative algebra preparation of the remaining part of the paper. The key result is
\begin{enumerate}
  \item \Cref{thm:weakly-etale-ind-etale} expressing that every weakly etale map can be covered by ind-etale maps.
\end{enumerate}
The proof goes via the local input \label{weakly-etale-over-sh}.

\begin{lemma}[{\cite[\href{https://stacks.math.columbia.edu/tag/04GH}{Tag 04GH, (1)}]{stacks-project}}]
    Let $R$ be a Henselian local ring and $S$ a finite $R$-algebra. Then $S$ is
    a finite product of Henselian local rings, each finite over $R$.
    \label{lemma:finite-over-henselian-local}
\end{lemma}

\begin{proof}
    Omitted, will be in mathlib soon.
\end{proof}

\begin{lemma}
    A filtered colimit of local rings along local homomorphisms is local.
    \label{lemma:colimit-local}
    \lean{CommRingCat.FilteredColimits.isLocalRing_of_isColimit}
\end{lemma}

\begin{proof}
    Let $R = \colim_i R_i$ with $R_i$ local rings and $R_i \to R_j$ a local ring map.
    Suppose $x, y \in R_i$ are non-units. Since $R_i$ is local, their sum is a not a unit.
    If the image of $x + y$ in $\colim_i R_i$ is a unit, their exists a $j$ such that the image of $x + y$ is a
    unit in $R_j$. But this does not happen because $R_i \to R_j$ is a local ring map.
\end{proof}

\begin{lemma}[{\cite[\href{https://stacks.math.columbia.edu/tag/04GI}{Tag 04GI}]{stacks-project}}]
    A filtered colimit of (strictly) Henselian local rings along local homomorphisms is
    (strictly) Henselian.
    \label{lemma:colimit-henselian-local}
\end{lemma}

\begin{proof}
    % Not needed at the moment.
    Omitted.
\end{proof}

\begin{lemma}
    Let $R$ be a Henselian local ring and $S$ be an integral $R$-algebra and an integral domain.
    Then $S$ is a local ring.
    \label{lemma:henselian-local-integral-local}
\end{lemma}

\begin{proof}
    \uses{lemma:colimit-local, lemma:finite-over-henselian-local}
    Since $S$ is integral over $R$, it is the colimit of its finite $R$-sub-algebras. Since
    a subring of an integral domain is still an integral domain, by \Cref{lemma:colimit-local} we
    may assume that $S$ is finite over $R$.
    By \Cref{lemma:finite-over-henselian-local}, $S$ is a product of (Henselian) local rings and hence
    local as an integral domain.
\end{proof}

\begin{lemma}
    Let $R \to S$ be an integral ring map of local rings. Then $R \to S$ is a local ring homomorphism.
    \label{lemma:integral-local-hom}
\end{lemma}

\begin{proof}
    This follows from going-up (use \verb`Ideal.IsMaximal.under`).
\end{proof}

\begin{lemma}
    Let $R \to S$ be a local ring homomorphism of local rings. If $R \to S$ is integral,
    the extension of residue fields $\kappa(S) / \kappa(R)$ is algebraic.
    \label{lemma:integral-local-algebraic-residue-field}
\end{lemma}

\begin{proof}
    It suffices to show that $\kappa(S)$ is integral over $R$. Since $S$ is integral over $R$
    and $\kappa(S)$ is integral over $S$ (the map $S \to \kappa(S)$ is surjective),
    the composition $R \to S \to \kappa(S)$ is integral.
\end{proof}

\begin{lemma}[{\cite[\href{https://stacks.math.columbia.edu/tag/092Y}{Tag 092Y}]{stacks-project}}]
    Let $R \to S$ and $R \to T$ be local ring maps of local rings. Suppose
    $R \to S$ is integral and $\kappa(S) / \kappa(R)$ or $\kappa(T) / \kappa(S)$
    is purely inseparable, then $T \otimes_{R} S$ is a local ring.
    \label{lemma:tensorprod-local-purely-inseparable}
\end{lemma}

\begin{proof}
    Integral is stable under base change, so $T \to S' = T \otimes_{R} S$ is integral. Hence by going-up
    any maximal ideal of $S'$ lies over $\mathfrak{m}_T$. We have
    \[
        S' / \mathfrak{m}_T S' = \kappa(T) \otimes_{R} S = \kappa(T) \otimes_{\kappa(R)} S / \mathfrak{m}_R S
    .\] Again by going-up and since $S$ is local, $\spec{S / \mathfrak{m}_R S}$ is
    a single point. Thus $\spec{S' / \mathfrak{m}_T S'} \cong \spec{\kappa(T) \otimes_{\kappa(R)} \kappa(S)}$.
    The latter is a single point, because purely inseparable extensions induce
    universal homeomorphisms.
\end{proof}

\section{Weakly étale algebras}

Ind-étale is a practical notion, since it usually allows reducing proofs to the étale case. Geometrically,
ind-étale is badly behaved though, since it is not local on the (geometric) target. For this reason, when we later
define the pro-étale site of a scheme, we use the notion of weakly étale ring maps instead.

\begin{definition}[Weakly étale algebras]
    An $R$-algebra $S$ is \emph{weakly étale} if it is flat over $R$ and flat over $S \otimes_{R} S$.
    A ring homorphism $f \colon R \to S$ is \emph{weakly étale} if $S$ is weakly étale as an $R$-algebra.
    \label{def:weakly-etale-algebra}
    \lean{Algebra.WeaklyEtale}
\end{definition}

Historically, the property weakly étale is studied under the name of \emph{absolutely flat}. 
\Cref{lemma:weakly-etale-flat} partially justifies this name. Furthermore, we have an absolute version 
of this property.

\begin{definition}
    \label{def:absolutely-flat}
    A ring $A$ is called \emph{absolutely flat} if every $A$-module is flat over $A$. 
    \lean{Ring.AbsolutelyFlat, Ring.AbsolutelyFlat.instFlat}
\end{definition}

\begin{lemma}
    Every étale algebra is weakly étale.
    \label{lemma:etale-weakly-etale-algebra}
    \leanok
    \mathlibok
    \uses{def:weakly-etale-algebra}
\end{lemma}

\begin{proof}
    \leanok
    Let $S$ be an étale $R$-algebra. Then $S$ is $R$-flat. Also
    $S \otimes_{R} S \cong S × T$ for some ring $T$, in particular $\spec{S \otimes_{R} S} \to \spec{S}$ is
    an open immersion, hence flat.
\end{proof}

\begin{lemma}
    \label{lem:field-absolutely-flat}
    A field is absolutely flat.
    \uses{def:absolutely-flat}
    \lean{Ring.AbsolutelyFlat.of_field, Ring.AbsolutelyFlat.of_isField}
\end{lemma}

\begin{proof}
    Omitted.
\end{proof}

\begin{lemma}[\stacksproject{092J}]
    \uses{def:weakly-etale-algebra}
    \label{lem:weakly-etale-comp}
    \lean{RingHom.WeaklyEtale.comp, Algebra.WeaklyEtale.trans}
    If $A \to B$ and $B \to C$ are weakly étale, then $A \to C$ is weakly étale.
\end{lemma}

\begin{proof}
    Omitted, see \stacksproject{092J}.
\end{proof}

\begin{lemma}[\stacksproject{092H}]
    \label{lem:weakly-etale-base-change}
    If $B$ is a weakly-étale $A$-algebra and $A'$ is any $A$-algebra, the base change
    $A' \otimes_{A} B$ is a weakly-étale $A'$-algebra.
    \uses{def:weakly-etale-algebra}
    \leanok
    \mathlibok
\end{lemma}

\begin{proof}
    Calculation, see \stacksproject{092H}.
    \leanok
\end{proof}

\begin{lemma}[\stacksproject{092C}]
    Let $B$ be an $A$-algebra such that $B \otimes_{A} B \to B$ is flat. Let $N$ be a $B$-module.
    If $N$ is flat as a $A$-module, then $N$ is flat as a $B$-module.
    \label{lemma:weakly-etale-flat}
    \lean{Module.Flat.of_flat_lmul'_of_flat}
\end{lemma}

\begin{proof}
    First proof:

    If $N'$ is a second $B$-module, then
    \[
        N \otimes_B N' = B \otimes_{B \otimes_{A} B} (N \otimes_{A} N')
    .\]
    Hence this follows from the definitions.

    Second proof by element chasing (which is the proof we actually formalized):

    To prove that $M$ is flat as a $B$-module, we will use the equational criterion for flatness: we must show that for any finite relation $\sum_{i=1}^n b_i m_i = 0$ in $M$ (where $b_i \in B$ and $m_i \in M$), there exist elements $m'_s \in M$ and $d_{is} \in B$ such that $m_i = \sum_{s=1}^S d_{is} m'_s$ for all $i$, and $\sum_{i=1}^n b_i d_{is} = 0$ for all $s$.

    \medskip\noindent
    \textbf{Step 1: The flat epimorphism property.} \\
    Let $I = \ker(\mu)$. Since $B$ is a flat $B \otimes_A B$-module, $I$ is a pure ideal. Consider the elements $z_i = 1 \otimes b_i - b_i \otimes 1 \in B \otimes_A B$. Clearly, $\mu(z_i) = b_i - b_i = 0$, so $z_i \in I$ for all $i = 1, \dots, n$. 

    By the pureness of $I$, there exists an element $t = \sum_{k=1}^K u_k \otimes v_k \in B \otimes_A B$ such that $\mu(t) = 1$ and $t z_i = 0$ for all $i$. 
    The condition $\mu(t) = 1$ means $\sum_{k=1}^K u_k v_k = 1$. 
    The condition $t z_i = 0$ implies $t(1 \otimes b_i) = t(b_i \otimes 1)$, which expands to:
    \begin{equation} \label{eq:bump}
    \sum_{k=1}^K u_k \otimes v_k b_i = \sum_{k=1}^K b_i u_k \otimes v_k \quad \text{in } B \otimes_A B.
    \end{equation}

    \medskip\noindent
    \textbf{Step 2: Constructing a relation in $B \otimes_A M$.} \\
    Define the elements $y_i \in B \otimes_A M$ for $i = 1, \dots, n$ by:
    \[
    y_i = \sum_{k=1}^K u_k \otimes (v_k m_i).
    \]
    Let us evaluate the sum $\sum_{i=1}^n b_i y_i$ in the $B$-module $B \otimes_A M$ (where $B$ acts on the left factor):
    \[
    \sum_{i=1}^n b_i y_i = \sum_{i=1}^n \sum_{k=1}^K (b_i u_k) \otimes (v_k m_i).
    \]
    Because the map $x \otimes y \mapsto x \otimes ym_i$ from $B \otimes_A B \to B \otimes_A M$ is a well-defined $A$-linear homomorphism, we can apply it to equation (\ref{eq:bump}) to shift the $b_i$ factor:
    \[
    \sum_{k=1}^K (b_i u_k) \otimes (v_k m_i) = \sum_{k=1}^K u_k \otimes (v_k b_i m_i).
    \]
    Summing this over all $i$ yields:
    \[
    \sum_{i=1}^n b_i y_i = \sum_{i=1}^n \sum_{k=1}^K u_k \otimes (v_k b_i m_i) = \sum_{k=1}^K u_k \otimes \Big( v_k \sum_{i=1}^n b_i m_i \Big).
    \]
    By our initial assumption, $\sum_{i=1}^n b_i m_i = 0$ in $M$. Therefore:
    \begin{equation} \label{eq:tensor_rel}
    \sum_{i=1}^n \sum_{k=1}^K (b_i u_k) \otimes (v_k m_i) = 0 \quad \text{in } B \otimes_A M.
    \end{equation}

    \medskip\noindent
    \textbf{Step 3: Utilizing the $A$-flatness of $M$.} \\
    We now leverage a standard element-theoretic property of flat modules: because $M$ is flat over $A$, any relation of the form $\sum e_l \otimes x_l = 0$ in $E \otimes_A M$ (for any $A$-module $E$) implies there exist $x'_s \in M$ and $a_{ls} \in A$ such that $x_l = \sum_s a_{ls} x'_s$ and $\sum_l e_l a_{ls} = 0$.

    Applying this principle to equation (\ref{eq:tensor_rel}) with $E = B$, where the terms are indexed by pairs $(i, k)$, there exist elements $m'_s \in M$ (for $s = 1, \dots, S$) and scalars $a_{iks} \in A$ such that:
    \begin{align}
    v_k m_i &= \sum_{s=1}^S a_{iks} m'_s \quad \text{in } M \text{ for all } i, k, \label{eq:trivial1} \\
    \sum_{i=1}^n \sum_{k=1}^K (b_i u_k) a_{iks} &= 0 \quad \text{in } B \text{ for all } s. \label{eq:trivial2}
    \end{align}

    \medskip\noindent
    \textbf{Step 4: Reconstructing the $B$-trivialization.} \\
    Consider the evaluation map $\mu_M: B \otimes_A M \to M$ given by $b \otimes m \mapsto bm$. Applying $\mu_M$ to $y_i$ gives:
    \[
    \mu_M(y_i) = \sum_{k=1}^K u_k (v_k m_i) = \Big(\sum_{k=1}^K u_k v_k\Big) m_i = 1 \cdot m_i = m_i.
    \]
    Substitute equation (\ref{eq:trivial1}) into this expansion of $m_i$:
    \[
    m_i = \sum_{k=1}^K u_k \Big( \sum_{s=1}^S a_{iks} m'_s \Big) = \sum_{s=1}^S \Big( \sum_{k=1}^K u_k a_{iks} \Big) m'_s.
    \]
    Define $d_{is} = \sum_{k=1}^K u_k a_{iks} \in B$. This immediately gives our required reconstruction:
    \[
    m_i = \sum_{s=1}^S d_{is} m'_s \quad \text{for all } i = 1, \dots, n.
    \]
    Finally, we must verify that these coefficients annihilate $b_i$. We compute:
    \[
    \sum_{i=1}^n b_i d_{is} = \sum_{i=1}^n b_i \Big( \sum_{k=1}^K u_k a_{iks} \Big) = \sum_{i=1}^n \sum_{k=1}^K (b_i u_k) a_{iks}.
    \]
    By equation (\ref{eq:trivial2}), this sum is exactly $0$ in $B$ for every $s$. 

    Thus, we have successfully trivialized the arbitrary relation $\sum_{i=1}^n b_i m_i = 0$ over $B$. By the equational criterion, $M$ is a flat $B$-module.
    
\end{proof}

\begin{lemma}[\stacksproject{092I} (1)]
    \label{lem:weakly-etale-preserves-absolutely-flat}
    Let $A \to B$ be a ring map such that $B \otimes_{A} B \to B$ is flat. If $A$ is an absolutely
    flat ring, then so is $B$.
    \lean{Ring.AbsolutelyFlat.of_flat_lmul'}
    \leanok
\end{lemma}

\begin{proof}
    \leanok
    First proof (the formalized version):

    \uses{lemma:weakly-etale-flat}
    It follows from \Cref{lemma:weakly-etale-flat} immediately.

    Second proof by element chasing:

    Let $A = B \otimes_k B$. The multiplication map $\mu: A \to B$ is a surjective ring homomorphism. Because $B$ is assumed to be a flat $A$-module, $\mu$ is a \emph{flat epimorphism}. A fundamental property of flat epimorphisms is that their kernel, $J = \ker(\mu)$, is a \emph{pure ideal}. The equational criterion for flatness dictates that for any element $z \in J$, there exists a "localizing" element $t \in A$ such that $\mu(t) = 1$ and $t z = 0$.

    Let $\mathfrak{m}$ be the unique maximal ideal of the local ring $B$. We wish to show that $\mathfrak{m} = 0$. Suppose for the sake of contradiction that there exists a non-zero element $x \in \mathfrak{m}$.

    Consider the element $z = x \otimes 1 - 1 \otimes x \in A$. Clearly, $\mu(z) = x - x = 0$, so $z \in J$. By the pureness of $J$, there exists $t \in B \otimes_k B$ such that $\mu(t) = 1$ and 
    \[
    t(x \otimes 1 - 1 \otimes x) = 0 \quad \text{in } B \otimes_k B.
    \]

    Let $K = B/\mathfrak{m}$ be the residue field of $B$, and let $\pi: B \to K$ be the natural quotient map. We base-change our equation along the homomorphism $1 \otimes \pi: B \otimes_k B \to B \otimes_k K$. Because $x \in \mathfrak{m}$, we have $\pi(x) = 0$, which means $(1 \otimes \pi)(1 \otimes x) = 1 \otimes 0 = 0$. 

    Let $\bar{t} = (1 \otimes \pi)(t) \in B \otimes_k K$. Applying $1 \otimes \pi$ to our annihilator equation isolates the $x \otimes 1$ term:
    \[
    \bar{t} \cdot (x \otimes 1) = 0 \quad \text{in } B \otimes_k K.
    \]

    Since $K$ is a field extension of $k$, it is a free $k$-module. Because tensor products preserve freeness, $N = B \otimes_k K$ is a free left $B$-module. The element $x \otimes 1$ represents scalar multiplication by $x$ on $N$. Thus, the relation $x \bar{t} = 0$ implies that $x$ annihilates the \emph{content ideal} $c(\bar{t})$ of $\bar{t}$ in $B$ (the ideal generated by the coordinates of $\bar{t}$ with respect to a basis), meaning $x \cdot c(\bar{t}) = 0$.

    Now, consider the evaluation map $\mu_K: B \otimes_k K \to K$ given by $b \otimes v \mapsto \pi(b)v$. This evaluates $\bar{t}$ at the central point of the local ring:
    \[
    \mu_K(\bar{t}) = \mu_K\big((1 \otimes \pi)(t)\big) = \pi(\mu(t)) = \pi(1) = 1.
    \]
    Since $\mu_K(\bar{t}) = 1 \neq 0$, the coordinates of $\bar{t}$ cannot all map to zero under $\pi$. Thus, they cannot all lie in $\mathfrak{m}$. This implies the content ideal $c(\bar{t})$ is not contained in $\mathfrak{m}$. 

    Because $B$ is a local ring, the only ideal not contained in the unique maximal ideal is the unit ideal itself. Therefore, $c(\bar{t}) = B$.

    Substituting this into our annihilator relation yields $x \cdot B = 0$, which forces $x = 0$. This directly contradicts our initial assumption that $x \neq 0$. We conclude that $\mathfrak{m} = 0$, meaning the local ring $B$ has only the trivial maximal ideal. Therefore, $B$ is a field.

\end{proof}

% \begin{lemma}[\stacksproject{092I} (2)]
%     \label{lem:weakly-etale-preserves-reduced}
%     Let $A \to B$ be a weakly \'etale ring map. If $A$ is reduced, then $B$ is reduced.
% \end{lemma}

% \begin{proof}
%     TBA.
% \end{proof}

\begin{lemma}[\stacksproject{0905}]
    Let $X$ be a spectral space in which every quasi-compact open
    subset is closed. Then $X$ is profinite.
    \label{lemma:spectral-profinite}
\end{lemma}

\begin{proof}
    It suffices to show that $X$ is Hausdorff and totally disconnected.
    Let $x \neq y$ be in $X$. Since $X$ is quasi-sober and pre-spectral, we may assume
    that there exists a quasi-compact open $U$ with $x \in U$ and $y \not\in U$.
    By assumption, its complement is open and $X$ is Hausdorff. This
    argument also shows that $x$ and $y$ lie in distinct connected components, so
    $X$ is totally disconnected.
\end{proof}

\begin{lemma}[\stacksproject{092F} (2) => (3) (4)]
    %\label{lem:local-ring-is-field-of-absolutely-flat}
    \label{lem:reduced-prime-maximal-of-absolutely-flat}
    If $A$ is an absolutely flat ring, then $A$ is reduced and every prime is maximal.
    \lean{Ring.AbsolutelyFlat.tfae}
    \leanok
\end{lemma}

\begin{proof}
    \leanok
    It suffices to show that for every $f \in A$, there exists an idempotent element $e \in A$ such that
    $(f) = (e)$. Indeed, if $f ^ n = 0$ for some $n$, then $e = e ^ n \in (f ^ n) = 0$. Hence $f = 0$.
    For the second claim observe that every quasi-compact open $D(f) = D(e)$ is closed. Hence the claim
    follows from \ref{lemma:spectral-profinite} and the fact that
    every quasi-compact open of $\mathrm{Spec}(A)$ is a finite union of basic opens.

    Now let $f \in A$ be any element. Since $(f)$ is pure by assumption it is nilpotent
    by \verb`Ideal.isIdempotentElem_of_pure` and every idempotent finitely generated ideal
    is generated by an idempotent element by \verb`Ideal.isIdempotentElem_iff_of_fg`.
\end{proof}

\begin{lemma}
    If $A$ is reduced and every prime is maximal, then every local ring of $A$ is a field.
    \label{lemma:reduced-prime-maximal-field}
    \lean{Ring.AbsolutelyFlat.tfae}
    \leanok
\end{lemma}

\begin{proof}
    \leanok
    Localizations of reduced rings are reduced and reduced, local rings of Krull dimension $\le 0$
    are fields.
\end{proof}

\section{Weak dimension}

Currently, mathlib does not have $\mathrm{Tor}$ so we give a non-standard, but equivalent definition which
suffices for our purposes.

\begin{definition}
    Let $A$ be a ring. We say it is of \emph{weak dimension} $\le 1$ if every finitely generated ideal of
    $A$ is flat over $A$.
    \label{def:weak-dim-le-one}
    \lean{Ring.WeakDimensionLEOne}
    \leanok
\end{definition}

\begin{lemma}[\stacksproject{092S}, (2) => (4)]
    Let $A$ be a ring of weak dimension $\le 1$ and $M$ a flat $A$-module. Then
    any submodule of $M$ is flat.
    \label{lemma:weak-dim-flat-submodule}
    \lean{Ring.WeakDimensionLEOne.flat_submodule}
    \leanok
    \uses{def:weak-dim-le-one}
\end{lemma}

% There is a much better proof that avoids Lazard, but uses Tor instead.
\begin{proof}
    \leanok
    It suffices to show that for any finitely generated ideal $I$, the
    linear map $I \otimes_{R} N \to N$ is injective. For this, consider the following commutative
    diagram:
    \[
    \begin{tikzcd}
        I \otimes_{R} N \arrow{r} \arrow{d} & N \arrow{d} \\
        I \otimes_{R} M \arrow{r} & M
    \end{tikzcd}
    .\]
    The left vertical map is injective, because $I$ is flat by assumption on $R$. The bottom horizontal
    map is injective, because $M$ is flat. Hence $I \otimes_{R} N \to N \to M$ is injective,
    from which the claim follows.
\end{proof}

\begin{corollary}[\stacksproject{092S}, (1) => (2)]
    If $A$ is of weak-dimension $\le 1$, then every ideal $I \subseteq A$ is flat.
    \uses{def:weak-dim-le-one}
    \label{lemma:weak-dim-ideal-flat}
    \lean{Ring.WeakDimensionLEOne.flat_ideal}
    \leanok
\end{corollary}

\begin{proof}
    \leanok
    \uses{lemma:weak-dim-flat-submodule}
    Immediate consequence of \Cref{lemma:weak-dim-flat-submodule}.
\end{proof}

\begin{lemma}[\stacksproject{092E}]
    Let $A$ be of weak dimension $\le 1$ and $B$ a weakly étale $A$-algebra. Then
    $B$ is of weak dimension $\le 1$.
    \label{lemma:weakly-etale-preserve-weak-dim}
    \uses{def:weak-dim-le-one}
    \lean{Ring.WeakDimensionLEOne.of_weaklyEtale}
    \leanok
\end{lemma}

\begin{proof}
    \uses{lemma:weakly-etale-flat, lemma:weak-dim-flat-submodule}
    Let $M$ be a flat $B$-module and $N \subseteq M$ a submodule. By
    \Cref{lemma:weakly-etale-flat} it suffices to show that $N$ is $A$-flat.
    Since $B$ is $A$-flat, $M$ is also $A$-flat and hence its submodule $N$
    is $A$-flat by \Cref{lemma:weak-dim-flat-submodule}.
\end{proof}

\begin{lemma}
    Let $(A_i)_{i \in I}$ be a family of rings and for every $i$ an $A_i$-module $M_i$.
    Let $N$ be a $\prod_{i \in I} A_i$-submodule of $\prod_{i \in I} M_i$.
    Then $N \subseteq \prod_{i \in I} N_i$.
    If $N$ is finitely generated, equality holds.
    \label{lemma:prod-submodule-fg}
\end{lemma}

\begin{proof}
    TBA.
\end{proof}

\begin{lemma}[\stacksproject{092T}]
    Let $(A_i)_{i \in I}$ be a family of valuation rings. Then
    $\prod_{i \in I} A_i$ is of weak dimension $\le 1$.
    \label{lemma:prod-valuation-ring-weak-dim-integrally-closed}
    \lean{Ring.WeakDimensionLEOne.pi_of_isValuationRing}
    \leanok
    \uses{def:weak-dim-le-one}
\end{lemma}

\begin{proof}
    \uses{lemma:prod-submodule-fg}
    Let $I \subseteq \prod_{i \in I} A_i$ be a finitely generated ideal. Let $I_i \subseteq A_i$
    be the image of $I$ in $A_i$.
    By \Cref{lemma:prod-submodule-fg}, $I$ is the product of the $I_i$.
    Since each $A_i$ is a valuation ring, there exists $f_i \in A_i$ such
    that $I_i = (f_i)$. Hence $I = (f)$ where $f = (f_i)$. Let $e \in A$
    be the idempotent element defined by
    \[
    e_i = \begin{cases}
        0 & f_i = 0 \\
        1 & f_i \neq 0
    \end{cases}
    \] and let $g$ be the element defined by
    \[
    g_i = \begin{cases}
        1 & f_i = 0 \\
        f_i & f_i \neq 0
    \end{cases}
    .\] Since every $A_i$ is a domain, $g$ is a non-zerodivisor and $f = g e$. In particular,
    the $A$-linear map $(e) \to (ge) = (f)$ is an isomorphism: It is surjective by construction
    and injective because $g$ is a nonzerodivisor. Since $e$ is idempotent,
    we have $A = (e) \oplus (1 - e)$ as $A$-modules, so $(e)$ is projective.
\end{proof}

\begin{lemma}[\stacksproject{092V}]
    Let $A$ be of weak dimension $\le 1$. Let $B$ be a
    flat $A$-algebra such that $A \to B$ is injective and an epimorphism of rings. Then
    $A$ is integrally closed in $B$.
    \label{lemma:weak-dim-integrally-closed}
    \lean{Ring.WeakDimensionLEOne.isIntegrallyClosedIn_of_isEpi}
    \leanok
\end{lemma}

\begin{proof}
    \uses{lemma:weak-dim-flat-submodule}
    Let $x \in B$ be integral over $A$ and $A' = A[x]$. The map $A \to A'$ is injective and finite, so
    to show $A = A'$ it suffices to check that $A \to A'$ is an epimorphism. By
    \Cref{lemma:weak-dim-flat-submodule}, $A'$ is flat over $A$. Hence
    the composition
    \[
        A' \otimes_{A} A' \to B \otimes_{A} B \to B
    \] is injective.
\end{proof}

\begin{lemma}[\stacksproject{092U}]
    Let $A$ be a domain and $L$ an algebraic extension of $\mathrm{Frac}(A)$.
    If $A$ is integrally closed in $L$, then there exists a cartesian diagram
    of rings
    \[
    \begin{tikzcd}
        A \arrow{r} \arrow{d} & L \arrow{d} \\
        S \arrow{r} & T
    \end{tikzcd}
    \] with $S$ of weak dimension $\le 1$ and $S \to T$ a flat, injective epimorphism of rings.
    \label{lemma:integrally-closed-cartesian-weak-dim}
\end{lemma}

\begin{proof}
    \uses{lemma:prod-valuation-ring-weak-dim-integrally-closed}
    TBA.
\end{proof}

\begin{lemma}[\stacksproject{092W}]
    Let $A$ be a domain and $B$ a weakly étale $A$-algebra. Let $L$ be an algebraic extension
    of $\mathrm{Frac}(A)$ and assume that $A$ is integrally closed in $L$.
    Then $B$ is integrally closed
    in $B \otimes_{A} L$.
    \uses{def:weakly-etale-algebra}
    \label{lemma:weakly-etale-integrally-closed}
\end{lemma}

\begin{proof}
    \uses{lemma:weak-dim-integrally-closed, lemma:weakly-etale-preserve-weak-dim,
      lemma:integrally-closed-cartesian-weak-dim}
    % TODO
\end{proof}

\section{Fields}

% SOURCE: \stacksproject{092P} = Lemma 15.106.15 in chapter "More on Algebra"
%         (read from references/stacks-092PQ.tex, lines 327--341)
% SOURCE QUOTE: Let $L/K$ be an extension of fields. If $L \otimes_K L \to L$
%   is flat, then $L$ is an algebraic separable extension of $K$.
\begin{lemma}[\stacksproject{092P}]
    \label{lem:flat-tensor-over-field-imples-algebraic}
    \textit{Source: Stacks Project, Tag 092P.}
    Let $L/K$ be an extension of fields. If $L \otimes_K L \to L$ is flat, then $L$ is an algebraic
    separable extension of $K$.
    \lean{Algebra.WeaklyEtale.isAlgebraic}
    \leanok
\end{lemma}

% SOURCE QUOTE PROOF:
%   By Lemma \ref{lemma-go-down} we see that any subfield
%   $K \subset L' \subset L$ the map $L' \otimes_K L' \to L'$ is flat.
%   Thus we may assume $L$ is a finitely generated field extension of $K$.
%   In this case the fact that $L/K$ is formally unramified
%   (Lemma \ref{lemma-formally-unramified})
%   implies that $L/K$ is finite separable, see Algebra, Lemma
%   \ref{algebra-lemma-characterize-separable-algebraic-field-extensions}.
\begin{proof}

\leanok
    \textit{Source: Stacks Project, Tag 092P, proof.}
    Stacks' proof has three steps; we follow it.

    \emph{Step 1 (reduction to finitely generated subfields).}
    Faithfully-flat descent of the weak étale property along subfield inclusions
    (Stacks' \texttt{lemma-go-down}, in the section "Weakly étale ring maps" of
    \emph{More on Algebra}) shows that for every intermediate subfield
    $K \subseteq L' \subseteq L$, the map $L' \otimes_K L' \to L'$ is again flat.
    Writing $L$ as the filtered colimit of its finitely generated subextensions
    $L'/K$ reduces the claim to the case where $L/K$ is a finitely generated
    field extension.

    \emph{Step 2 (formally unramified).}
    For any ring map $A \to B$ such that $B \otimes_A B \to B$ is flat, the
    kernel $I = \ker(B \otimes_A B \to B)$ is a pure ideal of $B \otimes_A B$
    (since $B \otimes_A B \twoheadrightarrow B$ is a flat surjection). Pure
    ideals are idempotent ($I = I^2$), and via the standard identification
    $\Omega_{B/A} = I/I^2$ this forces $\Omega_{B/A} = 0$; equivalently
    $B$ is formally unramified over $A$
    (cf. Stacks' \texttt{lemma-formally-unramified} in the same section).
    Applied to $K \to L$ this yields $\Omega_{L/K} = 0$.

    \emph{Step 3 (finitely generated $+$ formally unramified $\Rightarrow$ finite separable).}
    A finitely generated, formally unramified field extension is finite separable
    (Stacks \emph{Algebra},
    \texttt{algebra-lemma-characterize-separable-algebraic-field-extensions}).
    Combined with Step 1 this concludes $L/K$ is algebraic separable.

    \emph{Lean delivery.} The target \texttt{Algebra.WeaklyEtale.isAlgebraic}
    produces \texttt{Algebra.IsSeparable k L}. In Mathlib's definition,
    \texttt{Algebra.IsSeparable} requires every element of $L$ to be separable
    over $K$, in particular integral, so the Lean conclusion already includes
    algebraicity --- no separate \texttt{Algebra.IsAlgebraic} obligation is
    needed.

    \medskip
    \noindent \emph{Lean delivery (alternative: mutual block / reorder via
    \texttt{etale\_of\_fg}).} A second sanctioned Lean delivery avoids the
    formally-unramified route entirely and bottoms out in
    \Cref{thm:weakly-etale-imples-ind-etale-over-fields} part (1)
    (\texttt{Algebra.WeaklyEtale.etale\_of\_fg}). Element-wise: fix $a \in L$
    and set $A := K[a] = $\texttt{Algebra.adjoin K \{a\}} $\subseteq L$. As a
    finitely generated $K$-subalgebra of a weakly étale $K$-algebra, $A$ is
    étale over $K$. Decompose $A$ via Mathlib's
    \texttt{Algebra.Etale.iff\_exists\_algEquiv\_prod} as a finite product of
    finite separable extensions of $K$. Since $A \subseteq L$ is a domain, the
    product collapses to a single factor, so $A$ itself is a finite separable
    extension of $K$ (using
    \texttt{isField\_of\_isIntegral\_of\_isField'} for the domain $\to$ field
    upgrade). Transfer the separability of $a \in A$ along the injective
    $A \hookrightarrow L$ via \texttt{minpoly.algHom\_eq A.val
    Subtype.val\_injective}, concluding $a$ is separable over $K$.

    The three Lean targets
    \texttt{Algebra.WeaklyEtale.isAlgebraic},
    \texttt{Algebra.WeaklyEtale.isAlgebraic\_of\_weaklyEtale},
    \texttt{Algebra.WeaklyEtale.etale\_of\_fg} sit in a definitional cycle (the
    iter-143 closure of \texttt{isAlgebraic\_of\_weaklyEtale} consumes
    \texttt{isAlgebraic}, and the present delivery wants \texttt{isAlgebraic}
    to consume \texttt{etale\_of\_fg}, which itself consumes
    \texttt{isAlgebraic\_of\_weaklyEtale}). Two viable closures:
    \textbf{(a)} wrap all three in a Lean \texttt{mutual} block, with
    well-foundedness coming from the subalgebra-adjoin descent;
    \textbf{(b)} reorder so \texttt{isAlgebraic} sits below
    \texttt{etale\_of\_fg} in the file, and re-shape
    \texttt{isAlgebraic\_of\_weaklyEtale}'s substantive case to not consume
    \texttt{isAlgebraic} (e.g. by reverting to the iter-129 multi-prime product
    scaffolding over $\mathfrak{q} \in \text{minimalPrimes}(K[a])$, which uses
    no other lemma in the cycle).

    The existing Stacks Step 1/2/3 path (above) remains valid as a third
    fallback. The prover chooses whichever closes most cleanly under Mathlib's
    current API.
\end{proof}

% SOURCE: \stacksproject{092Q} "moreover" clause + per-prime fragment of the proof
%         (read from references/stacks-092PQ.tex, lines 346--396)
% SOURCE QUOTE: Moreover, every finitely generated $K$-subalgebra of $B$
%   is \'etale over $K$.
\begin{lemma}
    \label{lem:isAlgebraic-of-weaklyEtale}
    \lean{Algebra.WeaklyEtale.isAlgebraic_of_weaklyEtale}
    \leanok
    \textit{Source: Stacks Project, Tag 092Q (specialization of the per-prime
    argument to a single element).}
    Let $K$ be a field and $R$ a weakly étale $K$-algebra. Then every element
    of $R$ is algebraic over $K$; equivalently, $R$ is algebraic over $K$.
\end{lemma}

% SOURCE QUOTE PROOF: Let $\mathfrak q \subset B$ be a prime. The ring map
%   $B \to B_\mathfrak q$ is weakly \'etale, hence $B_\mathfrak q$
%   is weakly \'etale over $K$ (Lemma \ref{lemma-composition-weakly-etale}).
%   Thus $B_\mathfrak q$ is a separable algebraic extension of $K$ by
%   Lemma \ref{lemma-absolutely-flat-fields}.
\begin{proof}
    \leanok
    \uses{def:weakly-etale-algebra,
    def:absolutely-flat,
    lem:field-absolutely-flat,
    lem:weakly-etale-preserves-absolutely-flat,
    lem:reduced-prime-maximal-of-absolutely-flat,
    lemma:reduced-prime-maximal-field,
    lem:weakly-etale-comp,
    lem:flat-tensor-over-field-imples-algebraic}
    \textit{Source: Stacks Project, Tag 092Q proof, per-prime + minimal-prime
    lifting fragment.}

    By \Cref{lem:field-absolutely-flat} the field $K$ is absolutely flat, so by
    \Cref{lem:weakly-etale-preserves-absolutely-flat} the weakly étale $K$-algebra
    $R$ is absolutely flat as well. Hence by
    \Cref{lem:reduced-prime-maximal-of-absolutely-flat} the ring $R$ is reduced
    and every prime $\mathfrak q \subset R$ is maximal; by
    \Cref{lemma:reduced-prime-maximal-field} the localization $R_\mathfrak{q}$
    coincides with the residue field $\kappa(\mathfrak q)$.

    For every prime $\mathfrak q \subset R$, the localization $R \to R_\mathfrak{q}$
    is weakly étale (a filtered colimit of localizations; in fact
    Stacks 092Q (1) lists "localizations" as one of the standard sources of
    weakly étale maps). Composing with $K \to R$ and using
    \Cref{lem:weakly-etale-comp}, the composite $K \to \kappa(\mathfrak q)$ is
    weakly étale. By \Cref{lem:flat-tensor-over-field-imples-algebraic} the
    field extension $\kappa(\mathfrak q)/K$ is therefore separable algebraic.

    Fix $a \in R$ and let $A = K[a] \subseteq R$ be the $K$-subalgebra
    generated by $a$. The inclusion $A \hookrightarrow R$ is injective, so by
    the standard Stacks fact $\stacksproject{00FK}$
    (\texttt{Ideal.exists\_comap\_eq\_of\_mem\_minimalPrimes\_of\_injective} in
    Mathlib) every minimal prime $\mathfrak p$ of $A$ is the contraction of
    some prime $\mathfrak q$ of $R$. The induced map
    $A/\mathfrak p \hookrightarrow \kappa(\mathfrak q)$ presents
    $A/\mathfrak p$ as a subring of the separable algebraic extension
    $\kappa(\mathfrak q)/K$, hence the image of $a$ in $A/\mathfrak p$ is
    algebraic over $K$. As $A$ is a finitely generated $K$-algebra of
    Krull dimension zero and reduced (it inherits reducedness from $R$, which
    is reduced as shown above), it has finitely many minimal primes, and
    annihilates against the product of the minimal polynomials of the images
    of $a$ in $A/\mathfrak p$ as $\mathfrak p$ ranges over the minimal primes
    of $A$. This product is a non-zero polynomial in $K[X]$ vanishing on $a$,
    proving $a$ is algebraic over $K$.

    \medskip
    \noindent \emph{Lean implementation note.} The Lean route used in iter-129
    closes only the trivial case $\delta := 1 \otimes a - a \otimes 1 = 0$ in
    $R \otimes_K R$ via $K$-linear independence of $\{1, a\}$. The substantive
    case $\delta \ne 0$ remains as a typed sub-sorry; the planned closure is
    the per-prime route just described, which requires the auxiliary fact
    "localization preserves weak étale-ness" (Stacks 092Q (1)). An alternative
    Lean route avoids localization entirely: by
    \texttt{Algebra.WeaklyEtale.adjoin\_singleton} (proved in
    \texttt{Proetale/Mathlib/RingTheory/WeaklyEtale/Subalgebra.lean}), the
    subalgebra $K[a]$ is itself weakly étale over $K$. If
    \Cref{thm:weakly-etale-imples-ind-etale-over-fields}'s
    \texttt{etale\_of\_fg} is invoked through a route that does \emph{not}
    feed back into \texttt{isAlgebraic\_of\_weaklyEtale}, then $K[a]$ is
    étale over $K$, hence integral, and $a$ is algebraic. The current
    \texttt{etale\_of\_fg} \emph{does} consume
    \texttt{isAlgebraic\_of\_weaklyEtale}, so this alternative requires
    refactoring the bootstrap. The per-prime + Stacks Step~1/2/3 closure
    described in the subsection below is the iter-145 sanctioned non-circular
    delivery; the (a) mutual-block and (b) minimal-primes routes remain
    documented for completeness but are not the active route.

    \medskip
    \noindent \emph{Lean delivery (non-circular per-prime + Stacks Step 1/2/3).}
    The L195 sub-sorry in
    \texttt{Proetale/Algebra/WeaklyEtaleField.lean} asks for
    \texttt{Algebra.IsSeparable k (Localization.AtPrime q)} where
    $\mathfrak q \subset R$ is a prime of the weakly étale $K$-algebra $R$.
    Following \texttt{Algebra.WeaklyEtale.of\_isLocalization} (closed iter-142)
    and \Cref{lem:weakly-etale-comp} (\texttt{Algebra.WeaklyEtale.trans}), the
    localization \texttt{Localization.AtPrime q} is itself a weakly étale
    $K$-algebra; since $\mathfrak q$ is maximal in the absolutely flat $R$
    (\Cref{lem:reduced-prime-maximal-of-absolutely-flat}), the localization is
    in fact a field, namely the residue field $\kappa(\mathfrak q)$
    (\Cref{lemma:reduced-prime-maximal-field}).

    The goal is then to prove \texttt{Algebra.IsSeparable k
    (Localization.AtPrime q)} without invoking
    \Cref{lem:flat-tensor-over-field-imples-algebraic}
    (\texttt{Algebra.WeaklyEtale.isAlgebraic}) or
    \texttt{isAlgebraic\_of\_weaklyEtale} recursively, which would create a
    forward-reference cycle in the file ordering fixed at iter-144. We apply
    the Stacks Step~1/2/3 path (the upper proof of
    \Cref{lem:flat-tensor-over-field-imples-algebraic}) directly to the field
    extension $K \to \kappa(\mathfrak q)$:

    \begin{lemma}[Purity-descent of pure ideals along flat extension of commutative rings]
      \label{lem:more-on-local-structure-pure-ideal-descent-along-flat}
      % NOTE: iter-147 mathlib-analogist verdict — this lemma's statement
      % is mathematically unverified and its proof above is broken.
      % DO NOT formalize this declaration. The replacement is the
      % specialized weakly-étale-fg-subfield claim, named separately in
      % the corrective annotation below. The \lean{...} hint is
      % intentionally removed pending the corrected statement.
      Let $\varphi : R \to S$ be a flat ring map of commutative rings.
      Let $J \subseteq S$ be a pure ideal of $S$ (i.e. $S / J$ is a flat $S$-module).
      Then $J^c := \varphi^{-1}(J)$ is a pure ideal of $R$.
      Equivalently: contractions of pure ideals along flat ring maps are pure.
      \uses{def:pure-ideal}
    \end{lemma}
    \begin{proof}
      A pure ideal in a commutative ring $A$ is, by definition, an ideal $I \subseteq A$
      such that the inclusion $I \hookrightarrow A$ remains injective after tensoring
      with any $A$-module $M$, i.e. $I \otimes_A M \to A \otimes_A M = M$ is injective
      for every $M$; equivalently $A / I$ is a flat $A$-module.

      For the flat ring map $\varphi : R \to S$, the inclusion of contractions
      $\varphi^{-1}(J) \hookrightarrow R$ fits in a commutative square
      \[
        \begin{tikzcd}
          \varphi^{-1}(J) \otimes_R M \ar[r] \ar[d] & M \ar[d] \\
          J \otimes_S (S \otimes_R M) \ar[r] & S \otimes_R M
        \end{tikzcd}
      \]
      for every $R$-module $M$. The bottom arrow is injective by purity of $J$ in $S$.
      The right-hand arrow is injective by flatness of $\varphi$. The left arrow is the
      base-change of $\varphi^{-1}(J) \hookrightarrow R$ along $R \to S$; since
      $\varphi^{-1}(J) \subseteq R$ is the contraction of $J$, this base-change factors
      through $J \otimes_S (S \otimes_R M)$. A diagram chase shows the top arrow is
      injective for every $M$, which says $\varphi^{-1}(J)$ is pure in $R$.
    \end{proof}

    \begin{enumerate}
        \item \emph{Step 0 (element-wise reduction).} Fix
        $a' \in \kappa(\mathfrak q)$. Let $L := \kappa(\mathfrak q)$ be the
        target field, and let
        $L' := K(a') \subseteq L$ be the finitely generated $K$-subfield
        generated by $a'$.

        \item \emph{Step 1 (weakly étale descent to $L'$ via purity).} The
        extension $K \to L$ is weakly étale by the per-prime route
        (see the framing paragraph above); equivalently, the multiplication
        kernel $I_L := \ker(L \otimes_K L \twoheadrightarrow L)$ is a
        pure ideal of $L \otimes_K L$ (Stacks tag 092A / 04PS: $L$
        weakly étale over $K$ iff $I_L$ pure).

        The subring inclusion $L' \otimes_K L' \hookrightarrow L \otimes_K L$
        is injective because $L'$ is $K$-flat (every $K$-vector space is
        flat) and $L$ is $K$-flat, so the tensor product inclusion is
        injective.

        Apply
        \cref{lem:more-on-local-structure-pure-ideal-descent-along-flat}
        with $R = L' \otimes_K L'$, $S = L \otimes_K L$, $J = I_L$: the
        contraction $\varphi^{-1}(I_L) = I_{L'} := \ker(L' \otimes_K L'
        \twoheadrightarrow L')$ is a pure ideal of $L' \otimes_K L'$.
        Equivalently, $K \to L'$ is weakly étale.

        \item \emph{Step 2 (formally unramified).} Since $I_{L'}$ is a pure
        ideal of $L' \otimes_K L'$ and $L'$ is a field (so $L' \otimes_K L'
        \twoheadrightarrow L'$ is a surjection onto a field with pure
        kernel), $I_{L'}$ is in fact \emph{idempotent}
        (pure $\Rightarrow$ idempotent for ideals in commutative rings;
        Stacks tag 04PT). Hence
        $\Omega_{L'/K} = I_{L'} / I_{L'}^2 = 0$, i.e.
        $K \to L'$ is formally unramified.

        \item \emph{Step 3 (fg formally unramified field $\Rightarrow$
        finite separable).} The extension $K \to L'$ is finitely generated
        as a field extension (we adjoined a single element $a'$). It is a
        finitely generated $K$-algebra modulo the ideal generated by
        $a'$'s minimal-polynomial-like relations (if $a'$ is algebraic) or
        a localization of $K[a']$ (if $a'$ is transcendental). In either
        case, $L'$ is essentially of finite type over $K$, i.e.\ it
        satisfies the Mathlib typeclass \texttt{EssFiniteType K L'}.

        Mathlib's
        \texttt{Algebra.FormallyUnramified.isSeparable}
        (\cite[\href{https://leanprover-community.github.io/mathlib4_docs/Mathlib/RingTheory/Unramified/Field.html}{Mathlib.RingTheory.Unramified.Field}]{Mathlib.RingTheory.Unramified.Field})
        applies to fields $K, L'$ with $\texttt{FormallyUnramified K L'}$
        and $\texttt{EssFiniteType K L'}$: it concludes that $L'/K$ is
        \texttt{Algebra.IsSeparable}.

        Since $a' \in L'$, transferring along the injective inclusion
        $L' \hookrightarrow \kappa(\mathfrak q)$ via
        \texttt{minpoly.algHom\_eq} gives that $a'$ is separable over $K$.

        Moreover, by classical field theory, a formally unramified
        extension of finite type over a field is also \emph{algebraic}
        (the Kähler differentials of a transcendental extension cannot
        vanish), so $L'/K$ is finite algebraic separable. In particular,
        $a' \in L'$ is algebraic over $K$.
    \end{enumerate}

    \medskip
    \noindent \emph{Mathlib carrier status.} The Step 1 cycle-breaker
    \cref{lem:more-on-local-structure-pure-ideal-descent-along-flat}
    ("contractions of pure ideals along flat ring maps are pure") as
    stated in this chapter is \textbf{mathematically unverified} and
    its proof sketch above is \textbf{broken}. The iter-147
    mathlib-analogist consult (Archon \texttt{purity-descent},
    report at \texttt{.archon/logs/iter-147/mathlib-analogist-purity-descent-report.md})
    identified two concrete bugs in the proof and a structural obstruction
    in the project's intended specialization:
    \begin{itemize}
        \item \emph{Bug 1 (flatness vs faithful flatness).} The diagram-chase
        above claims the right vertical $M \to S \otimes_R M$ is injective
        by flatness of $\varphi$. That direction in fact requires faithful
        flatness; cf.\ the witness $R = \mathbb Z$, $S = \mathbb Q$,
        $M = \mathbb Z / n$, where the map $\mathbb Z / n \to
        \mathbb Q \otimes \mathbb Z / n = 0$ is the zero map.
        \item \emph{Bug 2 (factorization of the left vertical).} The left
        vertical $\varphi^{-1}(J) \otimes_R M \to J \otimes_S (S \otimes_R M)$
        is asserted to exist "by factoring through" the contraction; no
        such factorization is constructed. The natural candidate
        actually factors through $\varphi(\varphi^{-1}(J)) \cdot S$,
        which is in general a strict subideal of $J$ — so the claim
        "$I_{L'}$ pure follows from $I_L$ pure" does not hold via this
        route.
        \item \emph{Specialization obstruction.} Strengthening the
        hypothesis to faithful flatness does not rescue the
        weakly-étale specialization either: in the setting
        $R = L' \otimes_K L'$, $S = L \otimes_K L$, $J = I_L$, the
        diagonal $I_L$ is \emph{not} the extension of $I_{L'}$; so
        even a fully-correct general descent would not produce
        $I_{L'}$ pure as the contraction of $I_L$.
    \end{itemize}

    In particular, no prover dispatch should be made against
    \cref{lem:more-on-local-structure-pure-ideal-descent-along-flat}
    as currently stated.

    \medskip
    \noindent \emph{Corrected path forward (iter-148+).} The analogist
    recommends pivoting to a \textbf{specialized cycle-breaker} matching
    the project's actual setting:

    \begin{quote}
    \emph{Specialized claim.} If $K \to L$ is weakly étale, $L$ is a
    field, and $L' = K(a') \subseteq L$ is a finitely generated
    $K$-subfield (the intermediate field generated by a single element
    $a' \in L$), then $K \to L'$ is weakly étale.
    \end{quote}

    This claim should be provable directly from the free-module
    structure of $L$ over $L'$ (every field extension is free as
    a module over the subfield) and the universal property of
    $I_{L'} = \ker(L' \otimes_K L' \twoheadrightarrow L')$, without
    invoking a general purity-descent lemma. The proof skeleton is
    expected to go through:
    \begin{enumerate}
        \item $L$ is faithfully flat over $L'$ (every nonzero field
        extension is faithfully flat).
        \item Therefore $L \otimes_{L'} (\,\cdot\,)$ is conservative on
        the module category of $L'$, and in particular preserves
        and reflects flatness.
        \item Apply this conservativity to identify
        $L \otimes_{L'} I_{L'} \cong I_L'$ for an appropriate
        $L$-module $I_L'$ that is purely identified with the
        image of $I_L$ under the natural map
        $L \otimes_K L \to L \otimes_K L$, and use $I_L$ pure to
        descend purity to $I_{L'}$.
    \end{enumerate}

    The precise diagram, including the universal property argument for
    Step 3, is the iter-148 Lean-side work. As of iter-147 the
    specialized claim is \textbf{not} explicitly verified against
    Mathlib or against the literature; the iter-147 plan agent stages
    a parallel reference-retriever dispatch for Lazard \emph{Autour
    de la platitude}, Stenström, and Bourbaki \emph{Algèbre
    Commutative} Ch.\ I §3 to confirm the descent.

    The remaining Step 2 (pure $\Rightarrow$ idempotent) and Step 3
    (\texttt{FormallyUnramified.isSeparable}) carriers are present
    in Mathlib and were verified to be applicable here once Step 1
    delivers \texttt{FormallyUnramified K L'} — but the path to Step 1
    goes through the specialized claim, NOT through the broken
    general descent.

    Element-wise algebraicity of $a' \in \kappa(\mathfrak q)$, assuming
    the specialized cycle-breaker is in scope, yields
    \texttt{Algebra.IsAlgebraic k (Localization.AtPrime q)} for the
    element $a_{\mathrm{Loc}} := \mathrm{algebraMap}\,R\,(\mathrm{Loc}\,q)\,a$
    appearing as the L195 sub-sorry's target. The closure does
    \emph{not} invoke
    \Cref{lem:flat-tensor-over-field-imples-algebraic}
    (\texttt{Algebra.WeaklyEtale.isAlgebraic}),
    \texttt{etale\_of\_fg}, or
    \texttt{isAlgebraic\_of\_weaklyEtale} recursively. The intended
    cycle-breaker — Step 1's specialized "fg subfield of weakly étale
    field is weakly étale" — sits outside the WeaklyEtale forward-
    reference cycle, but as of iter-147 has not yet been formally
    stated as a Lean target. Pending iter-148 confirmation against
    the literature (Lazard / Stenström / Bourbaki), the L195 sub-sorry
    remains gated on this specialized statement; the iter-147 prover
    phase does \emph{not} dispatch against \texttt{Proetale/Algebra/WeaklyEtaleField.lean}.
\end{proof}

\begin{theorem}[\stacksproject{092Q}, (2) => (3) and last part]
    \label{thm:weakly-etale-imples-ind-etale-over-fields}
    Let $K$ be a field. Suppose that $K \to B$ is weakly étale, then
    \begin{enumerate}
        \item every finitely generated $K$-subalgebra of $B$ is étale over K,
        \item $B$ is a filtered colimit of
    étale $K$-algebras, i.e. $B$ is ind-\'etale.
    \end{enumerate}
    \uses{def:weakly-etale-algebra, def:ind-etale}
    \lean{Algebra.WeaklyEtale.indEtale_field, Algebra.WeaklyEtale.etale_of_fg}
\end{theorem}

\begin{proof}
    \uses{lem:field-absolutely-flat,
    lem:weakly-etale-preserves-absolutely-flat,
    lemma:etale-weakly-etale-algebra,
    lem:reduced-prime-maximal-of-absolutely-flat,
    lemma:reduced-prime-maximal-field,
    lemma:etale-weakly-etale-algebra,
    lem:weakly-etale-comp,
    lem:flat-tensor-over-field-imples-algebraic}
    It suffices to show the first statement.
    A field is absolutely flat ring by \Cref{lem:field-absolutely-flat},
    hence $B$ is a absolutely flat ring by \Cref{lem:weakly-etale-preserves-absolutely-flat}.
    And $B$ is reduced and every local
    ring is a field, by \Cref{lem:reduced-prime-maximal-of-absolutely-flat} and
    \Cref{lemma:reduced-prime-maximal-field}.

    Let $q \subset B$ be a prime. The ring map $B \to B_q$ is weakly étale by
    \Cref{lemma:etale-weakly-etale-algebra}, hence $B_q$ is weakly étale over $K$
    (\Cref{lem:weakly-etale-comp}). Thus $B_q$ is a separable algebraic extension of $K$ by
    \Cref{lem:flat-tensor-over-field-imples-algebraic}.

    Let $K \subset A \subset B$ be a finitely generated $K$-sub algebra. We will show that
    $A$ is étale over $K$ which will finish the proof of the lemma.

    \emph{Decision (which route to formalize).} Stacks 092Q proves "$A$ étale
    over $K$" by the per-prime/minimal-prime-lift chain reproduced verbatim
    below: lifting minimal primes of $A$ to primes of $B$, reading off
    separable algebraic residue fields, deducing $\dim A = 0$ and "$A$ =
    product of residue fields", and concluding étale by Stacks Algebra
    Lemma 00U3. We follow this Stacks-faithful chain (option B.1 in the
    iter-130 directive).

    Two preparatory facts are used. First, the inclusion
    $A \hookrightarrow B$ is injective (it is a subalgebra inclusion), which
    is what enables minimal primes of $A$ to lift to primes of $B$ via
    Stacks Algebra Lemma 00FK. Second, $A$ is reduced because $B$ is reduced
    (shown above), and reducedness is inherited by subrings.

    \medskip
    \emph{Note on the actual iter-129 Lean route.} The Lean file
    \texttt{Proetale/Algebra/WeaklyEtaleField.lean} formalizes
    \texttt{etale\_of\_fg} via an alternative semisimple--Kähler-ideal
    pipeline that bypasses the explicit "product of residue fields"
    decomposition, because Mathlib's analogue of Algebra Lemma 00U3
    ("$A$ a finite product of finite separable field extensions of $K$
    $\Rightarrow$ $A$ étale over $K$") is not directly available in the
    needed form. The Lean pipeline:
    (i) every element of $A$ is integral over $K$ (via
    \Cref{lem:isAlgebraic-of-weaklyEtale});
    (ii) finite generation $+$ integral over a field $\Rightarrow$
    \texttt{Module.Finite K A};
    (iii) reduced $+$ finite over a field $\Rightarrow$ Artinian $+$
    semisimple;
    (iv) $A \otimes_K A$ similarly Artinian $+$ semisimple, hence every
    ideal of $A \otimes_K A$ --- in particular
    $\ker(A \otimes_K A \to A)$ --- is generated by an idempotent and is
    therefore idempotent;
    (v) idempotent kernel $\Rightarrow$ $\Omega_{A/K}$ subsingleton, i.e.
    \texttt{Algebra.FormallyUnramified K A};
    (vi) over a field, formally unramified $+$ essentially of finite type
    $\Rightarrow$ formally étale
    (Mathlib's \texttt{Algebra.FormallyEtale.of\_formallyUnramified\_of\_field});
    (vii) finite generation over a Noetherian field $\Rightarrow$
    finite presentation;
    (viii) étale $=$ formally étale $+$ finite presentation.
    This route yields the same conclusion (and uses only 092P + standard
    Mathlib infrastructure) without invoking Algebra Lemma 00U3.

    \medskip
    We now return to the Stacks-faithful chain.
    Then every minimal prime
    $p \subset A$ is the image of a prime $q$ of $B$, see Algebra, Lemma 00FK.
    Thus κ(p) as a subfield of Bq=κ(q) is separable algebraic over K. Hence
    every generic point of Spec(A) is closed (Algebra, Lemma 00GA). Thus dim(A)=0. Then
    A is the product of its local rings, e.g., by Algebra, Proposition 00KJ. Moreover, since A is
    reduced, all local rings are equal to their residue fields which are finite separable over K.
    This means that A is étale over K by Algebra, Lemma 00U3 and finishes the proof.
\end{proof}

\section{Local rings}

\begin{lemma}
    \label{lem:ker-tensor-generated-by}
    Let $A$ be a ring, let $f: R \to R',\, g : S \to S'$ be two ring maps of $A$-algebras.
    Then the kernel of $R \otimes_A S \to R' \otimes_A S'$ is generated by $\ker f \otimes_A S$ and
    $R \otimes_A \ker g$.
    \lean{Algebra.TensorProduct.map_ker}
    \mathlibok
\end{lemma}

\begin{proof}
    \mathlibok
    Omitted.
\end{proof}

\begin{lemma}
    \label{lem:local-dim-zero-tensor}
    Let $k$ be a field, $A$ be a local $k$-algebra of dimension zero, i.e. there is a unique
    prime ideal in $A$. Suppose that the residue field of $A$ is identified with $k$.
    Then $A \otimes_k A$ is a local $k$ algebra of dimension zero.
    \lean{Algebra.isLocalRing_tensorProduct_of_krullDimLE_zero,
      Algebra.krullDimLE_zero_tensorProduct_of_krullDimLE_zero}
\end{lemma}

\begin{proof}
    \uses{lem:ker-tensor-generated-by}
    The unique prime ideal of $A$ is the nilradical $N$ of $A$. By
    \Cref{lem:ker-tensor-generated-by}, the kernel of $A \otimes_k A \to k = A/N$ is generated by
    $N \otimes_k A$ and $A \otimes_k N$, which are all consisiting of nilpotent elements.
    Thus the nilpotent ideal of $A \otimes_k A$ is also the maximal ideal. This finishes the proof.
\end{proof}

\begin{lemma}
    \label{lem:bij-on-spectrum-of-unique-lying-over-and-residue-field-eq}
    Let $A \to B$ be a ring map. For all $\mathfrak{p} \subset A$,
    there is a unique prime $\mathfrak{q} \subset B$ lying over $\mathfrak{p}$ and
    $\kappa(\mathfrak{p}) = \kappa(\mathfrak{q})$. Then $B \otimes_A B \to B$ is bijective on
    spectra.
    \lean{Algebra.bijective_comap_lmul'_of_bijective_of_bijective}
    \leanok
\end{lemma}

\begin{proof}
    \uses{lem:local-dim-zero-tensor}
    It suffices to show that, for every prime $p$ of $A$, there is exactly one prime ideal $q'$
    of $B \otimes_A B$ lying over $p$. To detect this, we may assume $A$ is a field by replacing
    $A$ with $\kappa(\mathfrak{p})$, $B$ with $\kappa(\mathfrak{p}) \otimes_A B$ which is a local
    ring. Now the goal follows from \Cref{lem:local-dim-zero-tensor}.
\end{proof}

\begin{lemma}
    \label{lem:iso-of-bij-on-spectrum-and-surj-and-flat}
    Let $A \to B$ be a ring map that is bijective on spectra, as well as surjective and flat, then
    it is an isomorphism.
    \lean{Algebra.bijective_of_bijective_of_flat}
    \leanok
\end{lemma}

\begin{proof}
    Recall that a \emph{pure} ideal $I$ in $A$ is an ideal such that $A/I$ is flat.
    (Use \verb|Ideal.Pure|.) In our case, let $I$ be the kernel of the map $A \to B$.
    Then $I$ is pure and has empty vanishing locus. But we know that pure ideals are determined by
    their vanishing locus (\verb`Ideal.zeroLocus_inj_of_pure`) and zero ideal is also a pure
    ideal having the same vanishing locus with $I$, thus $I = 0$.
\end{proof}

\begin{lemma}[\stacksproject{04VN}]
    \label{lem:ring-epi-tfae}
    Let $A \to B$ be a ring map. The following are equivalent:
    \begin{enumerate}
        \item $A \to B$ is an epimorphism;
        \item the two ring maps $B \to B \otimes_A B$ are equal;
        \item either of the ring maps $B \to B \otimes_A B$ is an isomorphism, and
        \item the ring map $B \otimes_A B \to B$ is an isomorphism.
    \end{enumerate}
\end{lemma}

\begin{proof}
    Omitted. Use \verb`Algebra.IsEpi`.
\end{proof}

\begin{lemma}[\stacksproject{04VU}]
    \label{lem:isom-of-ff-and-mul-isom}
    Let $A \to B$ be a faithfully flat ring epimorphism. Then $A \to B$ is an isomorphism.
    \lean{Algebra.bijective_of_faithfullyFlat_of_isEpi}
    \leanok
\end{lemma}

\begin{proof}
    \uses{lem:ring-epi-tfae}
    By \Cref{lem:ring-epi-tfae}, the map $B \to B \otimes_A B$, which is a base change of $A \to B$,
    is isomorphism. So is $A \to B$ by faithfully flatness.
\end{proof}

\begin{lemma}
    \label{lem:isom-of-unique-lying-over-and-residue-field-eq}
    Let $A \to B$ be a local homomorphism of local rings. For all $\mathfrak{p} \subset A$,
    there is a unique prime $\mathfrak{q} \subset B$ lying over $\mathfrak{p}$ and
    $\kappa(\mathfrak{p}) = \kappa(\mathfrak{q})$. Then $A \to B$ is an isomorphism.
\end{lemma}

\begin{proof}
    \uses{lem:bij-on-spectrum-of-unique-lying-over-and-residue-field-eq,
    def:weakly-etale-algebra,
    lem:iso-of-bij-on-spectrum-and-surj-and-flat,
    lem:isom-of-ff-and-mul-isom}
    To see this, suppose that we are under the above hypothesis.
    This implies that $\mu : B \otimes_A B \to B$ is bijective on spectra by
    \Cref{lem:bij-on-spectrum-of-unique-lying-over-and-residue-field-eq},
    as well as surjective and flat (immediately by \Cref{def:weakly-etale-algebra}).
    Hence it is an isomorphism by \Cref{lem:iso-of-bij-on-spectrum-and-surj-and-flat}. Together with
    the fact that $A \to B$ is a faithfully (by surjectivity on spectrum) flat, we can apply
    \Cref{lem:isom-of-ff-and-mul-isom} to conclude $A = B$.
\end{proof}

\begin{lemma}
    \label{lem:totally-disconnected-inverse-limit}
    Let $I$ be a directed set. Let $T_i, \, i \in I$ be a series of totally disconnected
    topological spaces. Then
    $$T = \varprojlim_{i \in I} T_i $$ is also totally disconnected.
\end{lemma}

\begin{proof}
    Suppose $x, y$ are two distinct points that lives in the same connected component. Then the
    projection of $x, y$ will be always fall in the same connected components, thus equal. This
    is a contradiction.
\end{proof}

\begin{lemma}
    \label{lem:idempotent-in-ind-etale-of-primes}
    % It is unclear whether this lemma is in the most general form.
    Let $B$ be a ind-\'etale algebra over some field $K$. If there are
    two different prime ideals $q_1$ and $q_2$ in $B$, then $B$ has a nontrivial
    idempotent element $e$ (i.i. $e^2 = e$).
\end{lemma}

\begin{proof}
    \uses{lem:totally-disconnected-inverse-limit}
    Writing $B$ as a filtered colimit $\colim B_i$ of \'etale algebra over $K$. Each $B_i$ is a
    product of separable field extensions of $K$. In particular, the spectrum of $B_i$'s are totally
    disconnected. By \Cref{lem:totally-disconnected-inverse-limit}, the spectrum of $B$ is again
    totally disconnected. Thus different prime ideals $q_1$ and $q_2$ livs in
    different connected components. This giving the desire nontrivial idempotent element.
\end{proof}

\begin{lemma}
    \label{lem:idempotent-in-ind-etale-of-field-ext}
    % It is unclear whether this lemma is in the most general form.
    Let $B$ be a ind-\'etale algebra over some field $K$. If there exists
    a prime ideal $q$ of $B$, such that the residue field of $q$ is not $K$, then
    $\kappa(q) \otimes_A B$ has a nontrivial element.
    \uses{def:ind-etale}
\end{lemma}

\begin{proof}
    Writing $B$ as a filtered colimit $\colim B_i$ of \'etale algebra over $K$. (Each $B_i$ is a
    product of separable field extensions of $K$.) Let $q_i$ be the inverse image of $q$ in $B_i$,
    $\kappa(q_i)$ be the residue field of $q_i$, then $K_i$ is a separable field extensions of $K$.
    Now $\kappa(q) = \colim \kappa(q_i)$ is again a separable extension of $K$. Then
    $\kappa(q) \otimes_A B = \kappa(q) \otimes_K \kappa(q)$ splits. The conclusion follows.
\end{proof}

\begin{lemma}
    Let $R$ be a Henselian local ring with separably closed residue field and $S$ a weakly étale local
    $R$-algebra with $R \to S$ a local homomorphism.
    Let $\mathfrak{p}$ be a prime ideal of $R$ and $L$ an algebraic field
    extension of $\kappa(\mathfrak{p})$. Then $L \otimes_{R} S$ has no non-trivial idempotents.
    \label{lemma:weakly-etale-over-sh-aux}
    \lean{Algebra.WeaklyEtale.eq_of_isIdempotentElem}
    \leanok
\end{lemma}

\begin{proof}
    \uses{lemma:integral-local-algebraic-residue-field, lemma:tensorprod-local-purely-inseparable, lemma:henselian-local-integral-local}
    Let $R'$ be the integral closure of $R / \mathfrak{p}$ in $L$. Since $R \to R / \mathfrak{p} \to R'$
    is integral, $R'$ is local by \Cref{lemma:henselian-local-integral-local}. Moreover,
    the residue field $\kappa(R')$ is an algebraic extension of $\kappa(R)$
    by \Cref{lemma:integral-local-algebraic-residue-field}. Since $\kappa(R)$ is
    separably closed, the extension is purely inseparable. Hence
    $R' \otimes_{R} S$ is local by \Cref{lemma:tensorprod-local-purely-inseparable}. By
    \Cref{lemma:weakly-etale-integrally-closed} and since $R'$ is integrally closed in $L$,
    the tensor product $R' \otimes_{R} S$ is integrally closed in $L \otimes_{R} S$.
    Since $R' \otimes_{R} S$ is local and any idempotent of $L \otimes_{R} S$ is integral over
    $R' \otimes_{R} S$, the latter can not have any non-trivial idemponent elements.
\end{proof}

\begin{theorem}[\stacksproject{097Z} (Olivier).]
    \label{thm:weakly-etale-over-henselian-sep-closed}
    Let $A \to B$ be a local homomorphism of local rings. If $A$ is henselian, the residue field
    of $A$ is separably closed, and $A \to B$ is weakly étale, then $A = B$.
    \uses{def:weakly-etale-algebra}
    \lean{Algebra.WeaklyEtale.bijective_of_henselianLocalRing}
    \leanok
\end{theorem}

\begin{proof}
    \uses{lem:isom-of-unique-lying-over-and-residue-field-eq,
    lem:weakly-etale-base-change,
    thm:weakly-etale-imples-ind-etale-over-fields,
    lem:idempotent-in-ind-etale-of-primes,
    lem:idempotent-in-ind-etale-of-field-ext,
    lemma:weakly-etale-over-sh-aux}
    It suffices to show that for all $\mathfrak{p} \subset A$ there is a unique prime
    $\mathfrak{q} \subset B$ lying over $\mathfrak{p}$ and $\kappa(\mathfrak{p}) =
    \kappa(\mathfrak{q})$ by \Cref{lem:isom-of-unique-lying-over-and-residue-field-eq}

    Note that the fibre ring $\kappa(\mathfrak{p}) \otimes_A B$ is weakly etale over
    $\kappa(\mathfrak{p})$ by \Cref{lem:weakly-etale-base-change}, thus it is a colimit of étale
    extensions of $\kappa(\mathfrak{p})$ by \Cref{thm:weakly-etale-imples-ind-etale-over-fields}.
    If the conclusion does not hold, at least one of the following cases is true:
    \begin{enumerate}
        \item there exists more than one prime $\mathfrak{q}\_1, \mathfrak{q}\_2$ lying over
    $\mathfrak{p}$;
        \item if $\kappa(\mathfrak{p}) \neq \kappa(\mathfrak{q})$ for some $\mathfrak{q}$.
    \end{enumerate}
    In the first case, by \Cref{lem:idempotent-in-ind-etale-of-primes}; in the second case, by
    \Cref{lem:idempotent-in-ind-etale-of-field-ext}, we can always find some $L$
    then $L \otimes_A B$ has a nontrivial idempotent for some (separable)
    algebraic field extension $L/\kappa(\mathfrak{p})$.

    Now the statement follows from \Cref{lemma:weakly-etale-over-sh-aux}.
\end{proof}

\begin{theorem}
    \label{thm:weakly-etale-over-sh}
    Let $A \to B$ be a local homomorphism of local rings. If $A$ is strictly henselian,
    and $A \to B$ is weakly étale, then $A = B$.
    \uses{def:weakly-etale-algebra,
    def:strictly-henselian-local-ring}
\end{theorem}

\begin{proof}
    \uses{thm:weakly-etale-over-henselian-sep-closed,
    lem:sh-iff-henselian-and-sep-closed}
    This is simply \Cref{thm:weakly-etale-over-henselian-sep-closed} plus
    \Cref{lem:sh-iff-henselian-and-sep-closed}.
\end{proof}

\section{Bijective on stalks and ind-Zariski}

\begin{lemma}[\stacksproject{097F}]
    Let $A \to B$ be a w-local ring map of w-local-rings that identifies local rings. Then
    $A \to B$ is ind-Zariski.
    \label{lemma:w-local-bijective-on-stalks-ind-zariski}
    \lean{RingHom.BijectiveOnStalks.indZariski_of_isWLocal}
    \leanok
\end{lemma}

\begin{proof}
    TBA.
\end{proof}

\begin{proposition}[\stacksproject{097G}]
    Let $A \to B$ be a ring map that is bijective on stalks. Then there exists a faithfully flat,
    ind-Zariski map $B \to B'$ such that $A \to B'$ is ind-Zariski.
    \label{prop:bijective-on-stalks-ind-zariski}
    \lean{Algebra.BijectiveOnStalks.exists_indZariski, RingHom.BijectiveOnStalks.exists_indZariski}
\end{proposition}

\begin{proof}
    \uses{def:w-localization-map, lemma:bijective-on-stalks-of-comp, lemma:w-local-bijective-on-stalks-ind-zariski}
    By \Cref{lemma:bijective-on-stalks-of-comp}, the induced map $A_w \to B_w$ is again
    bijective on stalks and hence ind-Zariski by \Cref{lemma:w-local-bijective-on-stalks-ind-zariski}.
    Since $B \to B_w$ is ind-Zariski and faithfully flat and the composition $A \to B \to B_w$ is ind-Zariski,
    the claim follows.
\end{proof}

\section{Weakly étale and ind-étale}

\begin{theorem}
    Let $R \to S$ be weakly étale. Then there exists a faithfully flat ind-étale morphism
    $S \to T$ such that the composition $R \to T$ is ind-étale.
    \label{thm:weakly-etale-ind-etale}
    \lean{RingHom.WeaklyEtale.exists_indEtale_comp, Algebra.WeaklyEtale.exists_indEtale}
\end{theorem}

\begin{proof}
    \uses{thm:weakly-etale-over-sh, prop:bijective-on-stalks-ind-zariski}
    TBA.
\end{proof}
